\subsection{Reconciliation Review}

\begin{longtable}{@{}L{0.25\textwidth}L{0.65\textwidth}@{}}
\caption{Reconciliation family summary}\\
\toprule
\textbf{Item} & \textbf{Details} \\
\midrule
\endfirsthead
\toprule
\textbf{Item} & \textbf{Details} \\
\midrule
\endhead
Purpose & Review mismatches between the gateway's internal records and external
provider evidence. \\
\bodyrowrule
Intended audience & Restricted / Internal operator. \\
\bodyrowrule
Authentication & Internal session cookie with operational access. \\
\bodyrowrule
Reference IDs & \apiid{REC-01}, \apiid{REC-02}. \\
\bottomrule
\end{longtable}

\apiblockheading{Endpoints.}
\begin{longtable}{@{}L{0.34\textwidth}L{0.56\textwidth}@{}}
\toprule
\textbf{Endpoint} & \textbf{Purpose} \\
\midrule
\endfirsthead
\toprule
\textbf{Endpoint} & \textbf{Purpose} \\
\midrule
\endhead
\code{GET /v1/ops/reconciliation} & List reconciliation records with optional
filters such as result, entity type, entity identifier, and limit. \\
\bodyrowrule
\code{GET /v1/ops/reconciliation/{record_id}} & Open a single reconciliation
record in detail. \\
\bodyrowrule
\code{POST /v1/ops/reconciliation/{record_id}/resolve} & Mark a reconciliation
record as resolved and record the review note. \\
\bottomrule
\end{longtable}

\apiblockheading{Business rules and usage notes.}
\begin{itemize}
  \item Reconciliation exists to preserve evidence when the gateway cannot
  safely auto-apply an external update.
  \item Resolution records the reviewer, the review note, and the final
  resolved state for later audit.
  \item Resolving a record acknowledges the review decision; it does not
  retroactively hide the original mismatch evidence.
\end{itemize}

\apiblockheading{Likely error codes.}
\begin{longtable}{@{}L{0.31\textwidth}L{0.59\textwidth}@{}}
\toprule
\textbf{Code} & \textbf{When it matters} \\
\midrule
\endfirsthead
\toprule
\textbf{Code} & \textbf{When it matters} \\
\midrule
\endhead
\code{RECONCILIATION_NOT_FOUND} & The requested reconciliation record does not
exist. \\
\bodyrowrule
\code{RECONCILIATION_ALREADY_RESOLVED} & The record has already been closed and
cannot be resolved again. \\
\bodyrowrule
\code{INTERNAL_AUTH_REQUIRED}, \code{INTERNAL_AUTH_INVALID_SESSION},
\code{INTERNAL_AUTH_FORBIDDEN} & The operator is not signed in or lacks the
required operational access. \\
\bottomrule
\end{longtable}

\apiblockheading{Happy-path example.}
\begin{lstlisting}[language=json,caption={Resolve reconciliation request example}]
{
  "actor": {
    "actor_type": "OPS",
    "actor_id": null,
    "reason": "Provider evidence reviewed."
  },
  "reviewed_by": null,
  "review_note": "Provider evidence reviewed and accepted."
}
\end{lstlisting}

\subsection{Manual Webhook Retry}

\begin{longtable}{@{}L{0.25\textwidth}L{0.65\textwidth}@{}}
\caption{Manual webhook retry family summary}\\
\toprule
\textbf{Item} & \textbf{Details} \\
\midrule
\endfirsthead
\toprule
\textbf{Item} & \textbf{Details} \\
\midrule
\endhead
Purpose & Retry a webhook event that has already exhausted or failed automatic
delivery. \\
\bodyrowrule
Intended audience & Restricted / Internal operator. \\
\bodyrowrule
Authentication & Internal session cookie with operational access. \\
\bodyrowrule
Reference IDs & \apiid{WH-10}. \\
\bottomrule
\end{longtable}

\apiblockheading{Endpoint.}
\begin{longtable}{@{}L{0.34\textwidth}L{0.56\textwidth}@{}}
\toprule
\textbf{Endpoint} & \textbf{Purpose} \\
\midrule
\endfirsthead
\toprule
\textbf{Endpoint} & \textbf{Purpose} \\
\midrule
\endhead
\code{POST /v1/ops/webhooks/{event_id}/retry} & Trigger a manual delivery retry
for a failed webhook event and record the operator's reason when supplied. \\
\bottomrule
\end{longtable}

\apiblockheading{Business rules and usage notes.}
\begin{itemize}
  \item Manual retry is allowed only when the webhook event is already marked as
  failed.
  \item A successful manual retry changes the delivery status to delivered.
  \item A failed manual retry keeps the event in failed state for further
  support action.
  \item Retrying a webhook does not change the underlying payment or refund
  outcome.
\end{itemize}

\apiblockheading{Likely error codes.}
\begin{longtable}{@{}L{0.31\textwidth}L{0.59\textwidth}@{}}
\toprule
\textbf{Code} & \textbf{When it matters} \\
\midrule
\endfirsthead
\toprule
\textbf{Code} & \textbf{When it matters} \\
\midrule
\endhead
\code{WEBHOOK_EVENT_NOT_FOUND} & The requested webhook event does not exist. \\
\bodyrowrule
\code{WEBHOOK_RETRY_NOT_ALLOWED} & The event is not currently in a failed state
and therefore is not eligible for manual retry. \\
\bodyrowrule
\code{INTERNAL_AUTH_REQUIRED}, \code{INTERNAL_AUTH_INVALID_SESSION},
\code{INTERNAL_AUTH_FORBIDDEN} & The operator is not signed in or lacks the
required operational access. \\
\bottomrule
\end{longtable}

\apiblockheading{Happy-path example.}
\begin{lstlisting}[language=json,caption={Manual webhook retry request example}]
{
  "actor_type": "OPS",
  "actor_id": null,
  "reason": "Retry after merchant endpoint recovered."
}
\end{lstlisting}

\begin{lstlisting}[language=json,caption={Manual webhook retry response example}]
{
  "event_id": "evt_01J...",
  "status": "DELIVERED",
  "attempt_count": 5,
  "last_attempt_result": "SUCCESS",
  "next_retry_at": null
}
\end{lstlisting}
